\section*{Abstract}

\policyengine{} UK is an open-source tax-benefit microsimulation model
covering \todo{count of modelled programs, e.g. from policyengine\_uk/programs.yaml (38 complete + 7 partial as of 2026-07-04)}
UK personal tax and benefit programs, built on the Family Resources Survey
(\frs{}) and validated against \ukmod{}, \hmrc{} and \dwp{} administrative
outturns, and the Survey of Personal Incomes (\spi{}). \euromod{}'s
Sutherland--Figari (2013) paper and \textsc{TAXSIM}'s Feenberg--Coutts (1993)
paper are the standard citations for those models; this paper is the
equivalent for \policyengine{} UK. It is \emph{versioned}: this version (v1)
documents the model as deployed in 2026, so that later engine migrations
produce a v2 rather than invalidating this description. We describe the
model's rules coverage and data construction (Section~\ref{sec:model},
Section~\ref{sec:data}), then validate calculated aggregates and
individual-level income tax liabilities against three independent sources:
a \ukmod{} cross-run and \hmrc{}/\dwp{}/\obr{} administrative and
official-forecast aggregates for major tax and benefit programs
(Section~\ref{sec:validation-ukmod}); \hmrc{}'s \spi{} 2020/21 individual
records, where \todo{restate the SPI match rate here once re-verified against
the current model version, citing the source notebook} of records match the
\spi{}'s reported income tax liability within \pounds{}10
(Section~\ref{sec:validation-spi}); and \frs{}-reported student loan
repayments against modelled repayments (Section~\ref{sec:validation-other}).
We additionally document two data-level validation records from the
\populace{} pipeline: the adjudicated comparison under which the current
default dataset replaced its predecessor
(Section~\ref{sec:validation-promotion}), and a fenced replay register
against \hmrc{}'s published personal-income distribution tables, in which
facts are asserted only where the survey instrument can express them
(Section~\ref{sec:validation-hmrc-replay}).
\todo{One-sentence summary of what the validation results collectively show,
once section~\ref{sec:validation} is written.}

\vspace{1em}
\noindent\textbf{Keywords:} microsimulation, tax-benefit model, United
Kingdom, model validation, survey microdata
